\documentclass[11pt,a6paper,french]{scrartcl}

\usepackage{babel}
\usepackage{microtype}

\usepackage{fontspec}
%\setmainfont[Ligatures=Discretionary]{EB Garamond}
%\setmainfont[Ligatures=Discretionary]{Palatino Linotype}
\setmainfont{EBGaramond}
[
  Extension      = .otf,
  UprightFont    = *-Regular,
  ItalicFont     = *-Italic,
  BoldFont       = *-Bold,
  BoldItalicFont = *-BoldItalic,
  Numbers        = Lowercase,
  Ligatures      = Discretionary,
  Style          = Swash
]
%\defaultfontfeatures{Ligatures=TeX}
%\setmainfont{UnifrakturMaguntia}
%\setmainfont{Junicode}
%\setmainfont{TeX Gyre Schola}
%\setmainfont{TeX Gyre Bonum}
%\setmainfont{TeX Gyre Termes}
%\setmainfont{Libertinus Serif}
%\setmainfont{Erewhon}
%\setmainfont{Vollkorn}
%\setmainfont{Cabin}
%\setmainfont{Medieval Sharp}
%\setmainfont{TeX Gyre Pagella}
%\setmainfont{Latin Modern Roman}
%\defaultfontfeatures{Ligatures=TeX}
%\setmainfont[Numbers=OldStyle]{Junicode}

\usepackage{verse}
%\usepackage{showframe}

\begin{document}
\thispagestyle{empty}
\settowidth{\versewidth}{Mon Dieu, quelle guerre cruelle !}

\poemtitle[]{\small%
\begin{tabular}{c}Plaintes d’un chrétien\end{tabular}\\
\begin{tabular}{c}sur les contrariétés\end{tabular}\\
\begin{tabular}{c}qu’il éprouve\end{tabular}\\
\begin{tabular}{c}au-dedans de lui-même\end{tabular}
}
\medskip
\begin{verse}[\versewidth]
Mon Dieu, quelle guerre cruelle !\\
Je trouve deux hommes en moi :\\
L’un veut que, plein d’amour pour toi,\\
Mon cœur te soit toujours fidèle ;\\
L’autre, à tes volontés rebelle,\\
Me révolte contre ta loi.

L’un, tout esprit et tout céleste,\\
Veut qu’au ciel sans cesse attaché,\\
Et des biens éternels touché,\\
Je compte pour rien tout le reste ;\\
Et l’autre, par son poids funeste,\\
Me tient vers la terre penché.

Hélas ! en guerre avec moi-même,\\
Où pourrai-je trouver la paix ?\\
Je veux, et n’accomplis jamais.\\
Je veux ; mais (ô misère extrême !)\\
Je ne fais pas le bien que j’aime,\\
Et je fais le mal que je hais.

Ô grâce, ô rayon salutaire !\\
Viens me mettre avec moi d’accord.\\
Et, domptant par un doux effort\\
Cet homme qui t’est si contraire,\\
Fais ton esclave volontaire\\
De cet esclave de la mort.
\end{verse}
\medskip
\begin{flushright}
Jean \textsc{Racine}, \textit{Cantiques spirituels}.
\end{flushright}

\end{document}
